% Übliche APA-7-Abschlussarbeit: muss im Image ohne tex_package_add kompilieren (Workflow „Dev image“, Smoke).
\documentclass[man]{apa7}
\usepackage[ngerman]{babel}
\usepackage{csquotes}
\usepackage[style=apa,backend=biber]{biblatex}
\addbibresource{literatur.bib}
\usepackage{booktabs}
\usepackage{threeparttable}
\usepackage{pgfplots}
\pgfplotsset{compat=1.18}
\usepackage{lineno}
\usepackage{float}
\usepackage{caption}
\usepackage{setspace}
\usepackage{hyperref}
\usepackage[ngerman]{cleveref}

\title{Lernzeit und Prüfungserfolg}
\shorttitle{Lernzeit}
\authorsnames{Alex Muster}
\authorsaffiliations{Universität Musterstadt}
\abstract{Die Arbeit prüft, ob mehr Lernzeit mit besseren Noten einhergeht.}
\keywords{Lernzeit, Prüfung}

\begin{document}
\maketitle
\linenumbers

\section{Einleitung}
Lernzeit gilt als wichtiger Faktor \parencite{muster2024}. \Cref{tab:noten} und \cref{abb:noten} fassen die Daten zusammen.

\begin{table}[H]
\caption{Mittlere Note nach Lernzeit}
\label{tab:noten}
\begin{threeparttable}
\begin{tabular}{lrr}
\toprule
Lernzeit & $n$ & Note \\
\midrule
unter 5 Stunden & 12 & 2,8 \\
5 Stunden und mehr & 15 & 2,1 \\
\bottomrule
\end{tabular}
\begin{tablenotes}
\item \textit{Anmerkung.} Noten von 1 (sehr gut) bis 5 (ungenügend).
\end{tablenotes}
\end{threeparttable}
\end{table}

\begin{figure}[H]
\caption{Mittlere Note nach Lernzeit}
\label{abb:noten}
\begin{tikzpicture}
\begin{axis}[ybar, symbolic x coords={unter 5 h, ab 5 h}, xtick=data, ymin=0, ylabel={Note}]
\addplot coordinates {(unter 5 h, 2.8) (ab 5 h, 2.1)};
\end{axis}
\end{tikzpicture}
\end{figure}

\printbibliography
\end{document}
